% =============================================================
% Conclusiones
% Apartado individual, no es un Capitulo
% Sugerido: resultados vs SAT, conclusion general, aporte a la ingenieria
% =============================================================
\chapter*{Conclusiones}
\label{cap:conclusiones}
\phantomsection
\addcontentsline{toc}{chapter}{Conclusiones}

\section*{Resultados obtenidos en relaci\'on a la Solicitud de Aprobaci\'on de Tema}

De los ocho objetivos espec\'ificos planteados en la Solicitud de
Aprobaci\'on de Tema (Secci\'on~\ref{sec:objetivos_especificos}), el
dise\~no, la implementaci\'on y la validaci\'on del circuito de
detecci\'on (primer y segundo objetivos) se encuentran completos
\cite{informeEneFeb2026,informeAbril2026}. La selecci\'on del conjunto
ADC/microcontrolador (tercer objetivo) tambi\'en se encuentra completa
\cite{informeEneFeb2026}.

PLACEHOLDER: completar esta secci\'on con el estado final de los
objetivos restantes (programaci\'on del microcontrolador, comunicaci\'on
con el seeder, software de decisi\'on, HMI en Python y validaci\'on con
mediciones del sistema LiDAR) una vez finalizado el proyecto.

\section*{Conclusi\'on general}

PLACEHOLDER: redactar la conclusi\'on general del trabajo una vez
finalizada la etapa de integraci\'on y validaci\'on descrita en el
Cap\'itulo~\ref{cap:modelo_experimental}.

\section*{Aporte a la ingenier\'ia}

PLACEHOLDER: como se se\~nala en la Secci\'on~\ref{sec:antecedentes}, no
se registran antecedentes de proyectos integradores de la FCEFyN
vinculados al desarrollo de sistemas de control para lazos de alta
resoluci\'on espectral de instrumentos LiDAR \cite{sat2026}; esta
secci\'on debe completarse explicitando el aporte concreto del sistema
desarrollado a este campo una vez finalizado el proyecto.
